\documentclass[10pt,a4paper]{amsart}
\usepackage[margin=19mm]{geometry}
\usepackage{amsmath,amssymb,mathtools}
\usepackage[scr=rsfs,cal=euler]{mathalfa}
\usepackage{xfrac}
\usepackage[unicode,hidelinks,bookmarks=false]{hyperref}
\newcommand{\ie}{i.e.\ }
\newcommand{\cf}{cf.\ }
\newcommand{\resp}{respectively\ }
\newcommand{\Subsection}{\subsection}
\providecommand{\nobreakdash}{\nobreak\hbox{-}}
\setlength{\emergencystretch}{2em}
\multlinegap0pt
\usepackage{bm}
\usepackage{bbm}
\usepackage{dsfont}
\usepackage{xspace}
\usepackage{cancel}
\usepackage{mathtools}
\usepackage[shortlabels]{enumitem}
\usepackage{amsthm}
\makeatletter
\@ifundefined{theorem}{\newtheorem{theorem}{Theorem}}{}
\makeatother
\title[An abstract criterion on the existence and global stability ]{An abstract criterion on the existence and global stability of stationary solutions for random dynamical systems and its applications}
\author{Xiang Lv}
\date{Source: arXiv:2508.08497v1 (CC BY 4.0)}
\begin{document}
\maketitle

\noindent\emph{Source and licence.} An abstract criterion on the existence and global stability of stationary solutions for random dynamical systems and its applications. Version arXiv:2508.08497v1 (CC BY 4.0), distributed under the Creative Commons Attribution 4.0 International licence, as recorded in the arXiv licence metadata for this exact version and shown on the abstract page https://arxiv.org/abs/2508.08497v1. Full rights notice: licenses/arxiv-2508.08497-CC-BY-4.0.txt.\par\smallskip

\noindent\textit{Editorial scope.} Counted unit: the Theorem whose source label is \texttt{thm3} in the article (source0.tex lines 445--450, source label \texttt{thm3}), delivered together with the article's own complete original proof of it (source0.tex lines 451--517, one \texttt{proof} environment of 67 lines). No mathematics is written, repaired or re-derived: statement and proof are the article's own text line for line. Required context retained uncounted: thm1 (lines 170--196); SDE2-1 (lines 394--396). Every definition, display and label of those lines is retained unchanged. Omitted from this package and not redistributed: 1--169, 197--393, 397--444, 518--869. Nothing inside a retained excerpt is omitted or altered: every retained body is delivered exactly as the article's lines are written, so no omission receipt of any kind accompanies this package. Citations: the retained excerpts carry no citation command at all, so no bibliography entry of the article is transcribed and no reference is redistributed. Reference adaptation: none was needed and none was made; every reference inside the delivered unit resolves inside it, so no unresolved reference or citation is produced. Printed slips kept literal: no printed word, symbol, hypothesis or number of the counted statement or of its proof was altered. Layout and class: the article's own class is \texttt{siamltex} with packages including amsmath, amssymb, amsfonts, color, xcolor, graphicx, manfnt, xy, enumerate, mathrsfs; the wrapper is the lane's pinned \texttt{amsart} editorial wrapper at 10pt on A4, which restates the theorem-like environments the retained text needs and adds the article's own macro definitions that the retained text needs (none), with extra wrapper packages (bm, bbm, dsfont, xspace, cancel, mathtools, enumitem). The delivered numbering is the unit's own. Assets: no retained span contains a figure environment, a graphics-inclusion command, a PDF-page inclusion, a table or a TikZ picture; no image file is copied into this package, the package declares no asset dependency and no image file is redistributed with it. Licence: version arXiv:2508.08497v1 of this article is distributed under the Creative Commons Attribution 4.0 International licence, as recorded in the arXiv licence metadata for this exact version and shown on the abstract page https://arxiv.org/abs/2508.08497v1. A full rights notice is delivered at licenses/arxiv-2508.08497-CC-BY-4.0.txt. Characters: every retained excerpt is pure ASCII, so the delivered engine, XeLaTeX with its default Latin Modern text font, reports no missing character.
\begin{theorem}\label{thm1} Let $(\theta,\varphi)$ be an RDS with the state space $X$.
Assume that $(\theta,\varphi)$ satisfies the following conditions:
\begin{enumerate}[{\rm(H1)}]
\item There exist a real number $\lambda>0$ and a random variable $\widehat{R}(\omega)$ such that for any $x,y\in X$,
\begin{equation}\label{H1}
d\bigl(\varphi(t,\theta_{-t}\omega,x),\varphi(t,\theta_{-t}\omega,y)\bigr)\leq\widehat{R}(\omega)e^{-\lambda t}d\bigl(x,y\bigr),\ t\geq t_0,\ \omega\in\Omega,
\end{equation}
where $t_0\geq0$ is a positive constant;
\end{enumerate}

\begin{enumerate}[{\rm(H2)}]
\item For any $x\in X$, there exists a $\lambda_0$-tempered random variable $R_x(\omega)$ such that
\begin{equation}\label{H2}
\sup_{t\geq0}d\bigl(\varphi(t,\theta_{-t}\omega,x),x\bigr)\leq R_x(\omega),\ \omega\in\Omega,
\end{equation}
where $0<\lambda_0<\lambda$.
\end{enumerate}
Then there exists a unique random equilibrium $U(\omega)$ of the RDS $(\theta,\varphi)$ such that
\begin{equation}\label{Conclu1}
\lim_{t\rightarrow\infty}\varphi(t,\theta_{-t}\omega,x)=U(\omega)\quad {\rm in}\quad X
\end{equation}
for any $x\in X$ and $\omega\in\Omega$. Moreover, the random equilibrium $U:\Omega\rightarrow X$ is tempered, i.e.,
\begin{equation}\label{Conclu2}
d\bigl(U(\omega),x\bigr)\leq R_x(\omega)
\end{equation}
for any $x\in X$ and $\omega\in\Omega$.
\end{theorem}

\begin{equation}\label{SDE2-1}
  dx(t)=g\bigl(x(t)\bigr)dt+\Sigma dB(t),
\end{equation}

\begin{theorem}\label{thm3} Assume that (A3) and (A4) hold, then there exists a unique random equilibrium $V_2(\omega)$ of the RDS $(\theta,\varphi_2)$ such that
\begin{equation}\label{eq30}
\lim_{t\rightarrow\infty}\varphi_2(t,\theta_{-t}\omega,x)=V_2(\omega)
\end{equation}
for any $x\in\mathbb{R}^n$ and $\omega\in\Omega$. Moreover, the random equilibrium $V_2:\Omega\rightarrow\mathbb{R}^n$ is tempered.
\end{theorem}

\begin{proof}
The proof will be divided into two parts. By (\ref{SDE2-1}), it follows naturally that
\begin{equation}\label{eq37}
\varphi_2(t,\omega,x)-\varphi_2(t,\omega,y)=x-y+\int_0^t\Bigl(g\bigl(\varphi_2(s,\omega,x)\bigr)-g\bigl(\varphi_2(s,\omega,y)\bigr)\Bigr)ds
\end{equation}
which together with (A3) implies that
\begin{align}\label{eq38}
&\frac{d}{dt}|\varphi_2(t,\omega,x)-\varphi_2(t,\omega,y)|^2\nonumber\\
=&2\left\langle\varphi_2(t,\omega,x)-\varphi_2(t,\omega,y),g\bigl(\varphi_2(t,\omega,x)\bigr)-g\bigl(\varphi_2(t,\omega,y)\bigr)\right\rangle\nonumber\\
\leq&-2L|\varphi_2(t,\omega,x)-\varphi_2(t,\omega,y)|^2.
\end{align}
Combining this and the Gronwall inequality, it is obvious that
\begin{equation}\label{eq39}
|\varphi_2(t,\theta_{-t}\omega,x)-\varphi_2(t,\theta_{-t}\omega,y)|\leq e^{-L t}|x-y|
\end{equation}
for all $x,y\in\mathbb{R}^n$, $t\geq0$ and $\omega\in\Omega$, which shows (H1).

Moreover, define
\begin{equation}\label{eq31}
z_2(t,\omega)\equiv z_2(\theta_t\omega)=\int_{-\infty}^te^{-(t-s)}\Sigma dB(s),
\end{equation}
which is the stationary solution (Ornstein-Uhlenbeck process) of the following affine SDEs
\begin{equation}\label{OU3}
dz_2(t)=-z_2(t)dt+\Sigma dB(t).
\end{equation}
Using It\^{o}'s formula, we can easily have that
\begin{equation}\label{eq32}
\varphi_2(t,\omega,x)-z_2(t,\omega)=x-z_2(\omega)+\int_0^t\Bigl(g\bigl(\varphi_2(s,\omega,x)\bigr)+z_2(s,\omega)\Bigr)ds,
\end{equation}
which yields that
\begin{align}\label{eq33}
&\frac{d}{dt}|\varphi_2(t,\omega,x)-z_2(t,\omega)|^2\nonumber\\
=&2\left\langle\varphi_2(t,\omega,x)-z_2(t,\omega),g\bigl(\varphi_2(t,\omega,x)\bigr)+z_2(t,\omega)\right\rangle\nonumber\\
=&2\left\langle\varphi_2(t,\omega,x)-z_2(t,\omega),g\bigl(\varphi_2(t,\omega,x)\bigr)-g\bigl(z_2(t,\omega)\bigr)\right\rangle\nonumber\\
&+2\left\langle\varphi_2(t,\omega,x)-z_2(t,\omega),g\bigl(z_2(t,\omega)\bigr)+z_2(t,\omega)\right\rangle\nonumber\\
\leq&-2L|\varphi_2(t,\omega,x)-z_2(t,\omega)|^2+L|\varphi_2(t,\omega,x)-z_2(t,\omega)|^2\nonumber\\
&+\frac1L|g\bigl(z_2(t,\omega)\bigr)+z_2(t,\omega)|^2\nonumber\\
=&-L|\varphi_2(t,\omega,x)-z_2(t,\omega)|^2+\frac1L|g\bigl(z_2(t,\omega)\bigr)+z_2(t,\omega)|^2.
\end{align}
Hence, by the Gronwall inequality and (A4), it is clear that
\begin{align}
&|\varphi_2(t,\omega,x)-z_2(t,\omega)|^2\nonumber\\
\leq& e^{-Lt}|x-z_2(0,\omega)|^2+\frac1L\int_0^te^{-L(t-s)}|g\bigl(z_2(s,\omega)\bigr)+z_2(s,\omega)|^2ds\nonumber\\
\leq& e^{-Lt}|x-z_2(\omega)|^2+\frac2L\int_0^te^{-L(t-s)}\left(|g\bigl(z_2(s,\omega)\bigr)|^2+|z_2(s,\omega)|^2\right)ds\nonumber\\
\leq& e^{-Lt}|x-z_2(\omega)|^2+\frac2L\int_0^te^{-L(t-s)}\left(2a^2|z_2(s,\omega)|^{2p}+2b^2+|z_2(s,\omega)|^2\right)ds\nonumber\\
\leq& e^{-Lt}|x-z_2(\omega)|^2+\frac2L\int_0^te^{-L(t-s)}\left(2a^2|z_2(s,\omega)|^{2p}+2b^2+|z_2(s,\omega)|^{2p}+1\right)ds\nonumber\\
\leq& \frac{4b^2+2}{L^2}+e^{-Lt}|x-z_2(\omega)|^2+\frac{4a^2+2}{L}\int_0^te^{-L(t-s)}|z_2(s,\omega)|^{2p}ds\nonumber
\end{align}
and then
\begin{align}\label{eq35}
&|\varphi_2(t,\theta_{-t}\omega,x)-z_2(t,\theta_{-t}\omega)|^2\nonumber\\
=&|\varphi_2(t,\theta_{-t}\omega,x)-z_2(\omega)|^2\nonumber\\
\leq&\frac{4b^2+2}{L^2}+e^{-Lt}|x-z_2(\theta_{-t}\omega)|^2+\frac{4a^2+2}{L}\int_0^te^{-L(t-s)}|z_2(s,\theta_{-t}\omega)|^{2p}ds\nonumber\\
=&\frac{4b^2+2}{L^2}+e^{-Lt}|x-z_2(\theta_{-t}\omega)|^2+\frac{4a^2+2}{L}\int_{-t}^0e^{Ls}|z_2(\theta_s\omega)|^{2p}ds\nonumber\\
\leq&\frac{4b^2+2}{L^2}+\sup_{t\geq0}\left\{e^{-Lt}|x-z_2(\theta_{-t}\omega)|^2\right\}+\frac{4a^2+2}{L}\int_{-\infty}^0e^{Ls}|z_2(\theta_s\omega)|^{2p}ds\nonumber\\
\triangleq&\widetilde R_x^2(\omega).
\end{align}
Note that $z_2$ is a tempered random variable and so also is $|z_2|^p$ for any $p\geq1$, we can directly verify that the random variable $\widetilde R_x^2$ is tempered.
Accordingly,
\begin{align}\label{eq36}
|\varphi_2(t,\theta_{-t}\omega,x)-x|^2
&\leq2|x|^2+4|\varphi_2(t,\theta_{-t}\omega,x)-z_2(\omega)|^2+4|z_2(\omega)|^2\nonumber\\
&\leq2|x|^2+4\widetilde R_x^2(\omega)+4|z_2(\omega)|^2\nonumber\\
&\triangleq R_x^2(\omega),\ t\geq0,\ x\in\mathbb{R}^n,
\end{align}
which together with Remark 1 gives (H2). Using Theorem \ref{thm1}, the proof is complete.
\end{proof}

\end{document}
